\documentclass[11pt]{article}
\usepackage[letterpaper,margin=0.88in]{geometry}
\usepackage{amsmath,amssymb,amsthm,mathtools}
\usepackage{microtype}
\usepackage{xcolor}
\usepackage[colorlinks=true,linkcolor=blue!55!black,urlcolor=blue!55!black,citecolor=blue!55!black]{hyperref}
\usepackage{enumitem}
\hypersetup{
  pdftitle={Banach-space well-separating common complements},
  pdfauthor={fa banach 001, agent lane 00, model GPT5.6}
}

\newtheorem{theorem}{Theorem}
\newtheorem{lemma}[theorem]{Lemma}
\newtheorem{corollary}[theorem]{Corollary}
\newtheorem{remark}[theorem]{Remark}
\newcommand{\R}{\mathbb R}
\newcommand{\N}{\mathbb N}
\newcommand{\dist}{\operatorname{dist}}

\title{A gliding affine-slab argument for\\
Banach-space well-separating common complements}
\author{Autonomous mathematical discovery run \texttt{fa\_banach\_001}\\
\texttt{agent\_lane\_00} --- model GPT5.6}
\date{13 August 2026}

\begin{document}
\maketitle

\begin{abstract}
Noethen asked whether every sequence of closed subspaces of the same finite
codimension in a real Banach space has a well-separating common complement,
and whether the tuples spanning such complements are prevalent.  We prove
both conjectures.  The missing codimension-one assertion follows from a
translated finite-dimensional slab estimate and a summable gliding
construction.  Quantitatively, for every sequence $(\varphi_j)$ of norm-one
functionals on an arbitrary Banach space there is a unit vector $x$ such that
$|\varphi_j(x)|\geq \exp[-O((\log(j+1))^2)]$ for every $j$.
\end{abstract}

\paragraph{Status.}
This is a candidate new full proof of Conjectures 3.4 and 3.15 in
\cite{Noethen}.  A bounded search found no later resolution; no priority claim
is made.

\section{The source questions}

Let $X$ be a real Banach space.  For closed subspaces
$V_j\subset X$ of common finite codimension $k$, a $k$-dimensional common
complement $C$ is \emph{well-separating} when there are $\delta_j>0$ such
that
\[
 \inf_{x\in C,\ \|x\|=1}\dist(x,V_j)\geq\delta_j,
 \qquad
 \lim_{j\to\infty}\frac{1}{j}\log\delta_j=0.
\]
Conjecture 3.4 of \cite{Noethen} asserts existence for every Banach space;
Conjecture 3.15 asserts prevalence of the tuples spanning such complements.
The source explicitly reduces existence to the Banach-space version of its
Lemma 3.6, concerning a sequence of norm-one functionals.  We prove a
quantitative strengthening of that lemma.

\section{A translated finite-batch estimate}

Write $B_2^n$ for the Euclidean unit ball and $\omega_n$ for its volume.

\begin{lemma}[affine-slab avoidance]\label{lem:finite}
Let $X$ be a real Banach space, let
$\varphi_1,\dots,\varphi_N\in X^*$ have norm one, and let $z\in X$ and
$r>0$.  There is $h\in X$ such that
\[
 \|h\|<r,
 \qquad
 |\varphi_j(z+h)|\geq \frac{r}{8N^2}
 \quad(1\leq j\leq N).
\]
\end{lemma}

\begin{proof}
Put $K=\bigcap_{j=1}^N\ker\varphi_j$ and let $q:X\to Y:=X/K$ be the
quotient map.  Then $n:=\dim Y\leq N$, and each $\varphi_j$ induces a
functional $\psi_j\in Y^*$ of norm one.  Indeed, this is the standard
isometric identification of $K^\perp$ with $(X/K)^*$.

By John's ellipsoid theorem, choose a linear isomorphism
$S:\R^n\to Y$ for which
\[
 S(B_2^n)\subset B_Y\subset \sqrt n\,S(B_2^n).
\]
Let $\ell_j=\psi_j\circ S$ and $a_j=\|\ell_j\|_2$.  Since
$\|\psi_j\|=1$ and $B_Y\subset\sqrt n\,S(B_2^n)$,
\begin{equation}\label{eq:aj}
 a_j\geq n^{-1/2}.
\end{equation}

Set $R=r/2$ and choose $U$ uniformly from $RB_2^n$.  For any real $b$,
an orthogonal change of coordinates and slicing perpendicular to $\ell_j$
give
\begin{equation}\label{eq:slab}
 \mathbb P\bigl(|\ell_j(U)+b|<\eta\bigr)
 \leq \frac{2\eta}{a_jR}\frac{\omega_{n-1}}{\omega_n}.
\end{equation}
The estimate is uniform in the translate $b$: the relevant interval has
length at most $2\eta/a_j$, and every ball section has volume at most
$\omega_{n-1}R^{n-1}$.  Also
$\omega_{n-1}/\omega_n\leq\sqrt n$.  For $n\geq2$, this follows by slicing
the unit ball and integrating over $|t|\leq n^{-1/2}$, since
$2(1-1/n)^{(n-1)/2}\geq1$; the case $n=1$ is immediate.
Together with \eqref{eq:aj}, \eqref{eq:slab} becomes
\[
 \mathbb P\bigl(|\ell_j(U)+b|<\eta\bigr)
 \leq \frac{4\eta n}{r}.
\]

Take $b=\varphi_j(z)$ and $\eta=r/(8N^2)$.  A union bound shows that the
probability that at least one of the $N$ inequalities fails is at most
\[
 N\frac{4\eta n}{r}=\frac{n}{2N}\leq\frac12.
\]
Hence some $u\in RB_2^n$ satisfies
$|\psi_j(qz+Su)|\geq\eta$ for all $j$.  The quotient vector $y=Su$ has
$\|y\|_Y\leq r/2$.  By the definition of the quotient norm it has a lift
$h\in X$ with $q(h)=y$ and
$\|h\|<\|y\|_Y+r/2\leq r$.  This lift has all the required properties.
\end{proof}

\section{One vector for the whole sequence}

\begin{theorem}[Banach version of Noethen's Lemma 3.6]\label{thm:sequence}
Let $X$ be an arbitrary real Banach space and let
$(\varphi_j)_{j\geq1}\subset X^*$ be norm-one functionals.  There are a unit
vector $u\in X$ and positive numbers $\delta_j$ such that
\[
 |\varphi_j(u)|\geq\delta_j\quad(j\geq1),
 \qquad
 \lim_{j\to\infty}\frac1j\log\delta_j=0.
\]
More precisely, one may take
\[
 \delta_j=\frac34,2^{-(m(j)^2+7m(j)-1)},
 \qquad
 m(j)=\max\{1,\lceil\log_2 j\rceil\}.
\]
\end{theorem}

\begin{proof}
Let $M_m=2^m$ and $\rho_1=1/4$.  Recursively set
\[
 \alpha_m=\frac{\rho_m}{8M_m^2},
 \qquad
 \rho_{m+1}=\frac{\alpha_m}{8}.
\]
Starting with $x_0=0$, apply Lemma~\ref{lem:finite} at stage $m$ to the
base point $x_{m-1}$, the first $M_m$ functionals, and radius $\rho_m$.
It supplies $h_m$ with $\|h_m\|<\rho_m$ such that, for
$x_m=x_{m-1}+h_m$,
\begin{equation}\label{eq:stage}
 |\varphi_j(x_m)|\geq\alpha_m
 \qquad(1\leq j\leq M_m).
\end{equation}

Since
\[
 \frac{\rho_{m+1}}{\rho_m}=\frac{1}{64M_m^2}\leq\frac1{256},
\]
the series $x=\sum_{m\geq1}h_m$ converges and
\[
 \|x\|<\sum_{m\geq1}\rho_m
 \leq \frac{1/4}{1-1/256}=\frac{64}{255}<1.
\]
Moreover, the tail after stage $m$ satisfies
\begin{equation}\label{eq:tail}
 \sum_{k>m}\rho_k
 \leq \frac{\rho_{m+1}}{1-1/256}
 =\frac{32}{255}\alpha_m
 <\frac14\alpha_m.
\end{equation}
Thus \eqref{eq:stage} and \eqref{eq:tail} imply
\begin{equation}\label{eq:limitbound}
 |\varphi_j(x)|>\frac34\alpha_m
 \qquad(1\leq j\leq M_m).
\end{equation}

The recurrence is explicit:
\[
 \alpha_1=2^{-7},
 \qquad
 \alpha_{m+1}=\alpha_m2^{-(2m+8)},
 \qquad
 \alpha_m=2^{-(m^2+7m-1)}.
\]
For $m=m(j)$, inequality \eqref{eq:limitbound} gives the stated $\delta_j$.
If $2^{m-1}<j\leq2^m$, then
$-\log\delta_j=O(m^2)=O((\log j)^2)$, so
$j^{-1}\log\delta_j\to0$.

Finally, \eqref{eq:limitbound} for $j=1$ shows that $x\neq0$.  Put
$u=x/\|x\|$.  Since $\|x\|<1$, normalization only increases all the
absolute functional values.
\end{proof}

\begin{remark}
No separability of $X$ is used: each correction lives over a
finite-dimensional quotient, and the corrections converge in $X$ itself.
For a complex Banach space, applying the theorem to the real norm-one
functionals $\operatorname{Re}\varphi_j$ also gives
$|\varphi_j(u)|\geq|\operatorname{Re}\varphi_j(u)|$, so the same conclusion
holds.
\end{remark}

\section{Resolution of both conjectures}

\begin{corollary}\label{cor:conjectures}
Conjectures 3.4 and 3.15 of \cite{Noethen} are true.  Namely, every sequence
of closed subspaces of a real Banach space having the same finite codimension
has a well-separating common complement, and the tuples spanning such
complements form a prevalent subset of the corresponding finite Cartesian
power of the space.
\end{corollary}

\begin{proof}
First suppose the subspaces $V_j$ are hyperplanes.  Choose
$\varphi_j\in X^*$ with $\|\varphi_j\|=1$ and
$V_j=\ker\varphi_j$.  For the unit vector $u$ from
Theorem~\ref{thm:sequence},
\[
 \dist(u,V_j)=|\varphi_j(u)|\geq\delta_j.
\]
Consequently $\operatorname{span}(u)$ is a well-separating common
complement.

For completeness, the arbitrary-codimension propagation in \cite{Noethen}
is purely Banach-space theoretic.  Inductively, for subspaces of codimension
$k+1$, enlarge each $V_j$ to a codimension-$k$ subspace and obtain a
well-separating $k$-dimensional complement $C_1$.  Then
$V_j\oplus C_1$ is a hyperplane, so the codimension-one result supplies a
common complementary line $C_2$.  Lemma 3.8 of the source, applied in
$X/V_j$, shows that $C_1\oplus C_2$ has a transversality lower bound
\[
 \frac{1}{2\sqrt5}\,\delta_j^{(1)}\delta_j^{(2)}.
\]
The product of two subexponentially decaying positive sequences is again
subexponential.  This proves Conjecture 3.4 for every finite codimension.

Finally, Proposition 3.20 of \cite{Noethen} proves in an arbitrary Banach
space that the existence of one well-separating common complement makes the
set of all spanning tuples prevalent.  Applying it to the complement just
constructed proves Conjecture 3.15.
\end{proof}

\section*{Verification notes}

The proof uses only John's ellipsoid theorem (already recorded as Theorem 2.5
in the source), elementary Euclidean slicing, quotient norms, and the two
Banach-space propagation results proved in the source.  The decisive feature
of Lemma~\ref{lem:finite} is uniformity in the affine translate $z$; this lets
successive finite-batch solutions be added without losing earlier margins.
The source conjecture pages and a separate constant-and-tail audit accompany
this packet.

\begin{thebibliography}{9}
\bibitem{Noethen}
F.~Noethen,
\emph{Well-separating common complements of a sequence of subspaces of the
same codimension in a Hilbert space are generic},
arXiv:1906.08514 (2019; PDF revision dated 9 March 2022).
\end{thebibliography}

\end{document}
